\section{\texorpdfstring{$G_2$}{G2} localization and local spectral operators}
\label{sec:g2-localization}
% Status: cited/open. G2 localization background is cited; local-kernel reductions remain open source calculations.

The \(G_2\) route supplies a local language for the DeVries program.  Its role
in this paper is local.  It gives a framework in which gauge fields,
charged matter, order-parameter data, and deformation modes can live at singular
loci of a compactification.  M-theory on \(G_2\)-holonomy spaces is a standard
route to four-dimensional \(\mathcal N=1\) structure; singularities provide the
gauge and matter sectors needed for particle physics
\cite{WittenG2Anomaly2001,AcharyaWitten2001,AcharyaGukov2004}.  The DeVries
question is whether a localized \(G_2\) operator can reduce to the same
two-channel determinant used in the pole-ratio clue.

\subsection{Singular support for gauge and matter sectors}
\label{sec:g2-singular-support}

The source-backed starting point is the localization of degrees of freedom.  Witten explains that \(G_2\) compactifications with singularities can support abelian or nonabelian gauge groups with chiral fermions, with anomaly inflow from the bulk supplying a consistency condition \cite{WittenG2Anomaly2001}.  Acharya and Witten describe examples where singular loci carry nonabelian gauge groups and conical or intersection points support chiral multiplets \cite{AcharyaWitten2001}.  Acharya and Gukov review the broader picture: singular \(G_2\) spaces are the essential setting for nonabelian gauge groups, chiral fermions, and strongly coupled gauge dynamics in M-theory compactifications \cite{AcharyaGukov2004}.

The DeVries problem asks for a local sector whose singular gauge locus,
localized charged matter or order-parameter mode, and singularity deformation
mode reduce to \(\mathcal M_J\).
The singular gauge locus is the natural address for adjoint/current data.  The localized charged matter or order-parameter mode is the natural address for the doublet side.  The deformation mode is the natural address for the scalar branch.

This local focus also controls the topology issue.  A compact \(G_2\) model imposes global constraints, including anomaly, representation, and line-operator data.  The determinant can be tested locally first.  A global compactification can then be asked to support the same local operator and the Standard Model global form.

\subsection{Higgs-bundle variables on an associative three-cycle}
\label{sec:g2-higgs-bundle-variables}

Modern local \(G_2\) gauge sectors can be organized by a Higgs-bundle-like system on an associative three-cycle \(M_3\).  Braun, Cizel, Hubner, and Schafer-Nameki describe a gauge field together with a one-form Higgs field on \(M_3\), with matter localized at critical loci and interactions controlled by gradient-flow data \cite{BraunEtAl2019}.  In the abelianized local description, the Higgs field can be written schematically as
\begin{equation}
  \Phi=\sum_{i=1}^n t_i\,df_i,
  \qquad
  \rho=\sum_{i=1}^n t_i\,\rho_i,
  \qquad
  \Delta f_i=\rho_i,
  \label{eq:g2-higgs-bundle-data}
\end{equation}
where \(t_i\) are Cartan generators and the \(\rho_i\) encode charge distributions.  For a representation with charge vector \(Q=(q_1,\ldots,q_n)\), the effective Morse function is
\begin{equation}
  f_Q=\sum_{i=1}^n q_i f_i,
  \qquad
  \rho_Q=\sum_{i=1}^n q_i\rho_i.
  \label{eq:g2-effective-morse}
\end{equation}
Matter in the representation \(R_Q\) localizes at points where
\begin{equation}
  df_Q(p)=0.
  \label{eq:g2-localization-condition}
\end{equation}
Flow lines between such critical points can generate mass terms, and trivalent
flow trees can generate Yukawa couplings when the representation charges sum to
zero \cite{BraunEtAl2019}.  The resulting local dictionary expresses gauge
data, matter localization, mass terms, and interactions as functions of the same
geometric input.

For the DeVries construction, the Higgs-bundle variables give a possible address for the local operator:
\begin{equation}
  \mathcal O_J^{G_2}
  =
  \mathcal O_J\big(W,\Phi,\rho_Q,f_Q\big),
  \label{eq:g2-operator-address}
\end{equation}
where \(W\) is the gauge field on \(M_3\).  The precise form of \(\mathcal O_J^{G_2}\) is open.  The variables \(W\), \(\Phi\), and \(f_Q\) already separate adjoint gauge data, scalar deformation data, and localized representation data.

The same Morse data suggest a Hodge/SUSY-QM refinement.  Witten's construction
uses \(d_t=e^{-th}de^{th}\), its adjoint, and the associated deformed Laplacian
to localize low modes near critical points of a Morse function
\cite{WittenMorse1982}.  In the local \(G_2\) setting, \(f_Q\) is already the
source-backed Morse function for matter localization.  A candidate DeVries
operator may therefore begin with a finite reduction of
\((\Omega^0,\Omega^1,d_t,d_t^\dagger)\) on the singular support.  The local
derivation must then supply the breaking entry, the two-channel projection, and
the electroweak assignment inside the same \(G_2\) model.

\subsection{Local two-channel kernel}
\label{sec:g2-two-channel-kernel}

Let \(a_J\) be an adjoint/current fluctuation supported along the singular gauge locus, and let \(h_J\) be a localized order-parameter fluctuation at a critical point of \(f_Q\).  Let
\begin{equation}
  \mathcal H_J^{G_2}=\operatorname{span}\{h_J,a_J\}.
  \label{eq:g2-light-sector}
\end{equation}
After integrating out heavier localized and bulk modes, the local inverse propagator or spectral kernel has the same abstract form as the endpoint and interval kernels:
\begin{equation}
  K_J^{G_2}(\lambda)=
  \begin{pmatrix}
  \lambda+\Sigma_{hh,J}^{G_2}(\lambda)
  &
  \Sigma_{ha,J}^{G_2}(\lambda)
  \\
  \Sigma_{ah,J}^{G_2}(\lambda)
  &
  \lambda+\Sigma_{aa,J}^{G_2}(\lambda)
  \end{pmatrix}.
  \label{eq:g2-local-kernel}
\end{equation}
The DeVries target is
\begin{equation}
  \Sigma_{hh,J}^{G_2}=0,\qquad
  \Sigma_{aa,J}^{G_2}=J,\qquad
  \Sigma_{ha,J}^{G_2}\Sigma_{ah,J}^{G_2}=J.
  \label{eq:g2-self-energy-target}
\end{equation}
The source-equation check for this target is explicit:
\begin{equation}
  \Phi=\sum_i t_i df_i,\qquad
  \rho=\sum_i t_i\rho_i,\qquad
  \Delta f_i=\rho_i,
  \label{eq:g2-source-equation-one}
\end{equation}
and
\begin{equation}
  f_Q=\sum_i q_i f_i,\qquad
  \rho_Q=\sum_i q_i\rho_i,\qquad
  df_Q(p)=0.
  \label{eq:g2-source-equation-two}
\end{equation}
The current entry \(\Sigma_{aa,J}^{G_2}=J\) is tested by the following local
ADE-pairing protocol.  Fix a source datum
\begin{equation}
  u_{G_2}^{\rm loc}
  =
  (M_3,\Gamma_{\rm ADE},W,\Phi,\rho_i,f_i,Q,p,\gamma,
  A_1\subset A_2,\Lambda_{G_2}),
  \label{eq:g2-local-source-datum}
\end{equation}
where the \(A_1\subset A_2\) enhancement records the local root-pairing test
\begin{equation}
  \operatorname{ad} A_2\!\downarrow_{A_1\times U(1)}
  =
  (\operatorname{ad} A_1)_0
  \oplus {\bf 1}_0
  \oplus {\bf 2}_{+q}
  \oplus {\bf 2}_{-q}.
  \label{eq:g2-a1-a2-enhancement-test}
\end{equation}
The singular support carries the pairing
\begin{widetext}
\begin{equation}
  \begin{aligned}
  \langle \eta,\zeta\rangle_{G_2,u}
  &=
  \int_{M_3}
  {\rm Tr}_{\rm ADE}\!\left(\eta\wedge *_{M_3}\zeta\right)
  +
  \sum_{p\in Z(df_Q)}
  \eta(p)^\dagger\,\mathsf G_Q(p)\,\zeta(p),
  \\
  Z(df_Q)&=\{p\in M_3\mid df_Q(p)=0\}.
  \end{aligned}
  \label{eq:g2-singular-support-pairing}
\end{equation}
\end{widetext}
The first term is the current-support pairing on the ADE locus.  The point
support term is the localized matter metric required by the later flow-overlap
entry.  The metric \(\mathsf G_Q(p)\) is theorem data: it must be positive on
the localized charge sector, gauge covariant, compatible with the charge vector
\(Q\), and reducible to the wavefunction or Hessian normalization at an isolated
Morse critical point.  A Morse--Bott critical component replaces the point
metric by the corresponding component metric and flow-moduli convention.

The local off-diagonal target is a flow-overlap theorem.  Let
\(\psi_{h,J}^{G_2}\) denote the localized order-parameter wavefunction and let
\(a_J^{G_2}\) denote the projected ADE current fluctuation.  Their source
normalizations are
\begin{equation}
  \langle\psi_{h,J}^{G_2},\psi_{h,J}^{G_2}\rangle_{G_Q}=1,
  \qquad
  \langle a_J^{G_2},a_J^{G_2}\rangle_{G_2,u}=1 .
  \label{eq:g2-flow-wavefunction-normalization}
\end{equation}
For an oriented flow line or flow-tree contribution \(\gamma\), introduce a
flow kernel \(\mathcal K_{\rm flow}^{G_2}(\gamma)\), orientation sign
\(\epsilon_\gamma\), and bilinear maps \(B_{ha}^{G_2}\) and \(B_{ah}^{G_2}\).
The hatted off-diagonal entries are theorem data:
\begin{equation}
  \begin{aligned}
  \widehat\Sigma_{ha,J}^{G_2}
  &=
  \Lambda_{G_2}^{-1}
  \sum_{\gamma}
  \epsilon_\gamma\,
  B_{ha}^{G_2}\!\left(
  \psi_{h,J}^{G_2},
  \mathcal K_{\rm flow}^{G_2}(\gamma)a_J^{G_2}
  \right),
  \\
  \widehat\Sigma_{ah,J}^{G_2}
  &=
  \Lambda_{G_2}^{-1}
  \sum_{\gamma}
  \epsilon_\gamma\,
  B_{ah}^{G_2}\!\left(
  a_J^{G_2},
  \mathcal K_{\rm flow}^{G_2}(\gamma)\psi_{h,J}^{G_2}
  \right).
  \end{aligned}
  \label{eq:g2-flow-overlap-definitions}
\end{equation}
The flow-overlap target is
\begin{equation}
  \widehat\Sigma_{ha,J}^{G_2}\widehat\Sigma_{ah,J}^{G_2}=J .
  \label{eq:g2-flow-overlap-target}
\end{equation}
The source datum is the existence of localized modes at \(df_Q=0\), oriented
flow-line mass terms, and trivalent flow-tree interaction data.  The theorem
specifies how the quadratic flow-line bilinear is extracted from those data.
The DeVries kernel then uses the product in
Eq.~\eqref{eq:g2-flow-overlap-target}.
The source-backed two-point object in Braun--Cizel--Hubner--Schafer-Nameki is
the Morse--Witten mass matrix
\begin{equation}
  M_{\rm MW}^{ab}
  =
  \sum_\gamma n_\gamma
  \exp[-tq(f(p_a)-f(p_b))],
  \qquad
  n_\gamma=\pm1,
  \label{eq:g2-braun-morse-witten-matrix}
\end{equation}
acting between localized chiral multiplet ground states.  Its Morse--Bott
extension keeps the same flow-line form for point-to-component trajectories and
gives a zero matrix element for two states localized on the same critical
circle.  The current source package therefore supplies a narrower test,
\begin{equation}
  \begin{gathered}
  M_{\rm MW}^{ab}
  \leadsto
  (M_{ha}^{\rm MW},M_{ah}^{\rm MW}),\\
  \mathcal H_J^{G_2}=\operatorname{span}\{h_J^{G_2},a_J^{G_2}\}.
  \end{gathered}
  \label{eq:g2-mw-to-ha-test}
\end{equation}
The source test isolates the failed derivation step: the Braun source data lack an
identification of the projected ADE current \(a_J^{G_2}\) as a
Morse--Witten endpoint state in this two-channel block.  The local flow-overlap
test applies only after an additional same-basis projection calculation supplies
\(M_{ha}^{\rm MW}\), \(M_{ah}^{\rm MW}\), \(\Lambda_{G_2}\), and the
orientation convention.

The local current test chooses a projected adjoint fluctuation
\(a_J^{G_2}\) with
\begin{equation}
  \langle a_J^{G_2},a_J^{G_2}\rangle_{G_2,u}=1
  \label{eq:g2-current-unit-norm}
\end{equation}
and defines
\begin{equation}
  \begin{aligned}
  \widehat K_J^{G_2,{\rm cur}}(\widehat\lambda)
  &=
  \Lambda_{G_2}^{-2}
  \left\langle a_J^{G_2},
  \mathcal K_{\rm cur}^{\rm ADE}(\lambda)a_J^{G_2}
  \right\rangle_{G_2,u},
  \\
  \lambda&=\Lambda_{G_2}^2\widehat\lambda .
  \end{aligned}
  \label{eq:g2-current-kernel-normalization}
\end{equation}
Here \(\mathcal K_{\rm cur}^{\rm ADE}\) denotes the projected quadratic
ADE-current operator induced on the singular gauge locus, including source
scale, projection, and possible singular-support boundary terms.  Its detailed
construction from the local \(G_2\) gauge theory is part of the source data.
The canonical current kernel is fixed by
\begin{equation}
  \begin{gathered}
  Z_J^{G_2}
  =
  \partial_{\widehat\lambda}
  \widehat K_J^{G_2,{\rm cur}}(0)>0,
  \\
  \widehat K_J^{G_2,{\rm cur,can}}
  =
  (Z_J^{G_2})^{-1}\widehat K_J^{G_2,{\rm cur}} .
  \end{gathered}
  \label{eq:g2-current-canonical-factor}
\end{equation}
The target is
\begin{equation}
  \widehat K_J^{G_2,{\rm cur,can}}(\widehat\lambda)
  =
  \widehat\lambda+\widehat\Sigma_{aa,J}^{G_2,{\rm can}}
  +O(\widehat\lambda^2),
  \qquad
  \widehat\Sigma_{aa,2}^{G_2,{\rm can}}=2 .
  \label{eq:g2-current-entry-target}
\end{equation}
This local test precedes the off-diagonal flow-overlap product and the
electroweak pole map.  Compact global completion, anomaly inflow, the charge
lattice, hypercharge embedding, and pole matching remain theorem data.
Appendix~\ref{app:target-g2-ade-pairing} states the same protocol as Target
IIIg.

\subsection{Dictionary for the \texorpdfstring{$G_2$}{G2} entries}
\label{sec:g2-entry-dictionary}

The local \(G_2\) dictionary is:
\begin{enumerate}
\item The current entry \(\Sigma_{aa,J}^{G_2}=J\) may come from an ADE singular gauge sector with an \(A_1\subset A_2\) enhancement and a singular-support pairing.  The source calculation must give \(\widehat\Sigma_{aa,2}^{G_2,{\rm can}}=2\).
\item The product \(\Sigma_{ha,J}^{G_2}\Sigma_{ah,J}^{G_2}=J\) may come from localized matter at \(df_Q=0\), a flow-line coupling, a bifundamental vev, or singularity-deformation mixing.  The source calculation must give the product from the point-support metric and local wavefunctions.
\item The reference entry \(\Sigma_{hh,J}^{G_2}=0\) may come from a critical point of \(f_Q\).  The source calculation must justify the zero entry through local normalization or a symmetry of the critical-point problem.
\item The ordered assignment \((J_H,J_{\rm adj})=(3/4,2)\) requires localized doublet/order-parameter data, adjoint/current data, photon-orthogonal completion data, hypercharge, a charge lattice, and a pole map in the same compact model.
\end{enumerate}
This dictionary turns \(G_2\) geometry into a determinant problem.  Each geometric
object must provide one entry of \(K_J^{G_2}\), one assignment rule, or one
consistency check.  A \(G_2\) route that supports the pole-ratio claim has to
derive at least one of \(\Sigma_{aa,J}^{G_2}\),
\(\Sigma_{ha,J}^{G_2}\Sigma_{ah,J}^{G_2}\), or \(\Sigma_{hh,J}^{G_2}\) from
Higgs-bundle or localized-singularity data with a stated inner product.
A \(G_2\) route enters the dimensional interpolation after the local
operator is embedded into a compact model whose chiral spectrum, anomaly
inflow, global form, and electroweak charge normalization are stated.

The normalization target includes the local inner product on the two-channel
space,
\begin{equation}
  \mathcal H_J^{G_2}
  =
  \operatorname{span}\{h_J^{G_2},a_J^{G_2}\},
  \qquad
  \langle h_J^{G_2},a_J^{G_2}\rangle_{G_2},
  \label{eq:g2-local-inner-product}
\end{equation}
where the pairing is the singular-support form of
Eq.~\eqref{eq:g2-singular-support-pairing}.  The current-entry derivation
identifies the gauge-field fluctuation \(W\) along the singular gauge locus and
then evaluates
\begin{equation}
  W\big|_{M_3}
  \quad\leadsto\quad
  \widehat K_J^{G_2,{\rm cur,can}}
  \quad\leadsto\quad
  \widehat\Sigma_{aa,2}^{G_2,{\rm can}}=2 .
  \label{eq:g2-W-current-entry}
\end{equation}
The rejection outputs are
\begin{equation}
  \widehat\Sigma_{aa,2}^{G_2,{\rm can}}\ne 2,\qquad
  Z_2^{G_2}\le0,\qquad
  \langle P_{a,2}^{G_2}a_2,P_{a,2}^{G_2}a_2\rangle_{G_2,u}\le0,
  \label{eq:g2-current-rejection-outputs}
\end{equation}
or an unresolved extra current channel.  These outputs classify the local ADE
pairing before a compactification or pole claim is made.
The off-diagonal-product derivation identifies the localized matter and
Higgs-bundle data,
\begin{equation}
  (\Phi,\rho_Q,f_Q,\psi_{h,J}^{G_2},\mathcal K_{\rm flow}^{G_2},\epsilon_\gamma)
  \quad\leadsto\quad
  \widehat\Sigma_{ha,J}^{G_2}\widehat\Sigma_{ah,J}^{G_2}=J,
  \label{eq:g2-higgsbundle-mixing-entry}
\end{equation}
through the bilinear definitions of
Eq.~\eqref{eq:g2-flow-overlap-definitions}.  Trivalent flow trees supply
interaction terms; the quadratic kernel requires the two-point flow-line
bilinear and its same-basis normalization.  The rejection outputs are
\begin{equation}
  \begin{gathered}
  \widehat\Sigma_{ha,J}^{G_2}\widehat\Sigma_{ah,J}^{G_2}\ne J,
  \qquad
  \mathsf G_Q\le0,\\
  \hbox{orientation convention absent},\\
  \hbox{extra coupled light channel}.
  \end{gathered}
  \label{eq:g2-flow-overlap-rejection}
\end{equation}
The order-parameter reference is derived from the
critical-point normalization of \(f_Q\),
\begin{equation}
  df_Q(p)=0
  \quad\leadsto\quad
  \Sigma_{hh,J}^{G_2}=0 .
  \label{eq:g2-critical-reference-entry}
\end{equation}
This is the local \(G_2\) analogue of the interval Schur-complement target:
the source theory must provide the three named entries before the singularity
language has derivational force.

\subsection{Higgsing as singular deformation}
\label{sec:g2-higgsing-deformation}

Acharya and Witten give a pattern directly relevant here: in examples where singular loci meet, gauge groups live on the singular components, chiral multiplets can be supported at the meeting point, and Higgsing corresponds to a geometric deformation or smoothing of the singularity \cite{AcharyaWitten2001}.  In their brane-related examples, a vev for a bifundamental field can break a product gauge symmetry to a diagonal subgroup, while the local singular geometry changes accordingly.

\begin{widetext}
For DeVries, this suggests the deformation chain
\begin{equation}
  \begin{aligned}
  \text{intersecting singular loci}
  &\longrightarrow
  \text{localized chiral/order-parameter mode}
  \\
  &\longrightarrow
  \text{deformed singularity}
  \longrightarrow
  K_J^{G_2}(\lambda).
  \end{aligned}
  \label{eq:g2-higgsing-chain}
\end{equation}
The scalar branch would then be read from the deformation direction.  If \(u_-(J)\) is the negative-branch eigenvector in \(\mathcal H_J^{G_2}\), the target scalar map is
\begin{equation}
  \begin{aligned}
  u_-(J)&\longmapsto \mathcal F_{G_2},
  \\
  \mathcal F_{G_2}
  &\in
  \left\{
  \begin{array}{l}
  \hbox{deformation parameter},\ f_Q(p),\ \rho_Q,\\
  \hbox{bifundamental vev},\ \hbox{localized modulus}
  \end{array}
  \right\}.
  \end{aligned}
  \label{eq:g2-negative-branch-map}
\end{equation}
The map must also respect the electroweak ray.  The deformation parameter can supply the scalar side, while the pole-ratio branch supplies the projective vector side.
\end{widetext}

\subsection{Anomaly and global checks}
\label{sec:g2-anomaly-global}

The \(G_2\) route has consistency tests.  Witten's anomaly-inflow discussion
links local singularity data to four-dimensional chiral matter and anomaly
cancellation \cite{WittenG2Anomaly2001}.  Therefore a local DeVries operator
has to pass the following checks:
\begin{enumerate}
\item the singular gauge locus supports the required adjoint sector;
\item localized charged matter has the desired representation content;
\item anomaly inflow or its local analogue is compatible with the charged fields;
\item the global gauge-group form of Sec.~\ref{sec:global-form} is compatible with the local representations;
\item flavor labels remain separate from compactification topology until a representation map is derived.
\end{enumerate}
These checks can reject a candidate completion even when the two-channel determinant is obtained locally.

\subsection{\texorpdfstring{$G_2$}{G2} route criterion}
\label{sec:g2-completion-criterion}

The \(G_2\) route has the same reduction chain as the endpoint and interval routes:
\begin{widetext}
\begin{equation}
\begin{aligned}
  \text{singular }G_2\text{ neighborhood}
  &\longrightarrow
  (W,\Phi,\rho_Q,f_Q)
  \longrightarrow
  K_J^{G_2}(\lambda),
  \\
  K_J^{G_2}(\lambda)
  &\longrightarrow
  \lambda^2+J\lambda-J
  \longrightarrow
  \{\sPole,\mathcal F_{G_2},\mathcal C_{\rm top}\}.
\end{aligned}
\label{eq:g2-completion-chain}
\end{equation}
\end{widetext}
The first arrow is source-backed by singular \(G_2\) compactification and Higgs-bundle literature.  The second arrow is the DeVries determinant derivation.  The final arrow is the pole-placement, scalar-branch, and global-form assignment.  The local determinant is the first compatibility condition; compact global data then supply anomaly, representation, and line-operator tests.
